% !TEX root = userguide.tex

\chapter{Constraints}
\pgst{tfemch}

\label{ch:constraints}

\def\chaptertitle{Constraints}
\def\headdate{18th February 2015}

In this chapter it is shown how to handle constraints in \tfem.

\section{Type of constraints}

Constraints need to be defined on some type of geometrical entity: 
a point, curve, surface, nodeset, elementset or object.
On such an entity
there are two types of constraints available in \tfem:
\emph{global} constraints and \emph{distributed} constraints\footnote{At the
time of writing
only global constraints
have been implemented 
on elementsets (see Sec.~\ref{sec:elementsets}).}.

Global
constraints act on all degrees of freedom in the geometrical entity and are
thus not tied to particular (nodal) points. For example imposing the flow rate
through a surface involves all degrees of freedom on that surface.

Distributed constraints are tied to particular points. For example connecting
two curves together requires a constraint in each point along the curves that
locally connects the degrees of freedom on both sides of the connection.

Both type of constraints (global and distributed) are not static and can
still contain unknowns. For example along a curve it is known that the degree
of freedom is a linear function of a coordinate along that curve, but the
function is unknown. This means that there are two \emph{additional unknowns}
(the two constants of the linear function) which are added to the system vector
of unknowns and solved along with all other degrees of freedom.

Examples:
\begin{enumerate}
\item Applying Dirichlet boundary conditions. In Fig.~\ref{fig:slip} a flow
geometry is given where along curves C1, C2 and C3 there is perfect slip along
the wall.
\begin{figure}[ht!]
   \begin{center}
      \includegraphics{slip}
   \end{center}
   \caption{Flow region with a curve not along the coordinate directions}
    \label{fig:slip}
\end{figure}
 The normal velocity component is zero:
\begin{equation}
    \vec u\cdot\vec n =0
  \label{eq:normalvel}
\end{equation}
Applying this condition along C1 and C3 is easy: essential boundary conditions
for the $u_y$ component. Along C2 a local coordinate
transformation can be used (see Chapter~\ref{ch:transformations}). An alternative is to specify Eq.~\eqref{eq:normalvel} as a constraint
along curve C2. This is a distributed constraint.

\item Periodic boundary conditions. Consider a flow region given in
Fig.~\ref{fig:channel}.
\begin{figure}[ht!]
   \begin{center}
      \includegraphics{channel}
   \end{center}
   \caption{Flow region with periodic boundary conditions}
    \label{fig:channel}
\end{figure}
If we assume periodic boundary conditions the
velocity profile along C1 is the same as C2:
\begin{equation}
    \vec u_1(y) = \vec u_2(y)
  \label{eq:periodic}
\end{equation}
This could be implemented by eliminating all unknowns on C2 by identifying them
with the corresponding unknowns on C1. This option is not available in \tfem.
Another option is to implement Eq.~\eqref{eq:periodic} as a constraint that
involves both C1 and C2. This is a distributed constraint.

\item Imposing a flow rate. It is sometimes more convenient to specify the flow
rate instead of the pressure difference. For example in Fig.~\ref{fig:channel}
we would like to impose:
\begin{equation}
     \int_{C1}u_x\D y = Q
  \label{eq:flowrate}
\end{equation}
where $Q$ is a constant. This is a global constraint on curve C1.

\item Unknowns in the constraint. In the driven cavity problem
Fig.~\ref{fig:driven2}
\begin{figure}[ht!]
   \begin{center}
      \includegraphics{driven_cavity2}
   \end{center}
   \caption{Driven cavity flow with unknown velocity $U$}
    \label{fig:driven2}
\end{figure}
the lid is driven by a known force. So we know that all velocities on the upper
curve must be the same $(U,0)$, but the value of $U$ is unknown. We can
implement this boundary condition using a constraint along the upper boundary
with an additional unknown $U$.
   
\end{enumerate}

\section{Implementation with Lagrange multipliers}

Constraints are implemented using \emph{Lagrange multipliers}. To illustrate
this, we consider the following minimization problem related to linear
elasticity problems (minimum potential energy): find the vector
$\col u_{\text{min}}$ that minimizes the function
\begin{equation}
   I(\col u) = \frac 12\col u^T\mat S\col u-\col u^T\col f
   \label{eq:I}
\end{equation}
where $\col u$ and $\col f$ are vectors of dimension $N$ and
 $\mat S$ is a $N\times N$ symmetric matrix.
We can find the minimum by variation (differentiation):
\[
 \begin{split}
    \delta I(\col u)&=\frac12\delta\col u^T\mat S\col u+\frac12\col u^T\mat S\delta\col u
    -\delta\col u^T\col f\\
    &=\delta\col u^T(\mat S\col u - \col f) = 0,\quad\text{for all }\delta \col u
 \end{split}
\]
or $\col u_{\text{min}}$ is the solution of the system of equations
\[
 \mat S\col u = \col f
\]
Now we constrain the minimization process with the following $N_c$ constraints
having $N_\text{add}$ additional unknowns:
\begin{equation}
    \mat A \col u + \mat B \col u_\text{add}=\col g
   \label{eq:constraint}
\end{equation}
where $\mat A$ is a $N_c\times N$ matrix, 
      $\mat B$ is a $N_c\times N_\text{add}$ matrix, $\col g$ is vector of
dimension $N_c$ and $\col u_\text{add}$ is a vector of additional unknowns 
of dimension $N_\text{add}$. The solution of this constrained minimization
process, i.e.\ the minimum of $I(\col u)$ in Eq.~\eqref{eq:I}
under the constraint Eq.~\eqref{eq:constraint} is given by the
\emph{stationary value} of the function (saddle point problem):
\[
   J(\col u,\col\lambda,\col u_\text{add})= I(\col u) +
    \col\lambda^T ( \mat A \col u + \mat B \col u_\text{add}-\col g )
\]
where (the components of) $\col\lambda$ are called \emph{Lagrange
multipliers}.

Sometimes there are external tractions/forces on the unknowns involved in the
constraint. In the formulation above the tractions/forces should be specified 
using the vector $\col f$. 
For example, on some curve the displacements are constrained to
be a linear functions given by an unknown position vector and a rotation. The
remaining external force and torque on this curve can be specified using
tractions on the curve and are added to $\col f$. However, in this problem it
would be more convenient to just specify the integral force and torque. This
can be done most easily by realizing that these integral quantities are just
the `dual' forces of the general displacements given by $\col u_\text{add}$. So
instead of adding tractions to $\col f$ we can equivalently find the
stationary value of the function
\[
  \begin{split}
   \hat{J}(\col u,\col\lambda,\col u_\text{add})&= 
        J(\col u,\col\lambda,\col u_\text{add})-\col u^T_\text{add}\col h \\
        &=\frac 12\col u^T\mat S\col u-\col u^T\col f
    +\col\lambda^T ( \mat A \col u + \mat B \col u_\text{add}-\col g )
                                               -\col u^T_\text{add}\col h
  \end{split}
\]
where $\col h$ is the vector of external forces on the additional unknowns.
The stationary value of $\hat J$ is given by
\[
   \delta \hat{J} = \delta \col u^T(\mat S\col u -\col f +\mat A^T\col\lambda)
  + \delta\col\lambda^T(\mat A\col u+\mat B\col u_\text{add}-\col g)
  + \delta\col u^T_\text{add}(\mat B^T\col\lambda-\col h) = 0
\]
for all $\delta\col u$, $\delta\col\lambda$ and $\delta\col u_\text{add}$,
which leads to the following system of equations
\begin{align*}
   \mat S\col u +\mat A^T\col\lambda &=\col f \\
\mat A\col u+\mat B\col u_\text{add}&=\col g\\
\mat B^T\col\lambda&=\col h
\end{align*}
or in matrix form
\begin{equation}
  \begin{bmatrix}
     \mat S & \mat A^T & \mat 0 \\
     \mat A & \mat 0 & \mat B   \\
     \mat 0 & \mat B^T & \mat 0
  \end{bmatrix}
  \begin{bmatrix}
     \col u \\
     \col\lambda \\
     \col u_\text{add}
  \end{bmatrix}
=
  \begin{bmatrix}
     \col f \\
     \col g \\
     \col h
  \end{bmatrix}
\label{eq:constraintsystem}
\end{equation}
Note that
\begin{itemize}
   \item The Lagrange multipliers add forcing to the original equations.
   \item The system is symmetric.
   \item The additional equations have zero blocks on the diagonal.
   \item The additional unknowns constrain the Lagrange multipliers. 
\end{itemize}

\section{Numerical implementation}

The matrices $\mat A$ and $\mat B$ and vectors $\col g$ and $\col h$ 
in Eq.~\eqref{eq:constraintsystem} need to be supplied to the system matrix by
an `element-routine' and
are assembled nodal-point-wise or element-wise, depending on the discretization
used (see below). Also the matrix $\mat A$ is split into two separate matrices
($\mat A_1$ and $\mat A_2$) if two
geometrical entities are
involved (like two curves, surfaces or objects that are connected) or
the constraint is on two-sided
objects (like in fluid-solid interaction problems):
\[
\mat A\col u = \mat A_1\col u_1+\mat A_2\col u_2
\]
where $\col u_1$ and $\col u_2$ are the unknowns on the first and second
entity or `side' of the object, respectively.

\subsection*{Global constraints}

Global constraints such as Eq.~\eqref{eq:flowrate} are usually build
element-wise along the curve, surface, elementset or
object\footnote{An object needs to
have a topology in order to be used in `element-wise' global constraints}.
The contribution to the integral is computed on element level and the matrix
$\mat A_e$ and vector $\col g_e$ need to be filled. The matrix $\mat A_e$ is of
size $N_c\times n_\text{fe}$, where $n_\text{fe}$ is the number of unknowns
in element $e$. The
right-hand side value $Q$ is set in $g_1$ (the first element) and $g_2$, $g_3$,
\ldots are set to zero.

Another possibility is to build global constraints nodal point
wise\footnote{This is the only possibility for a constraint in a point
or on a nodeset.}. For
example building the constraint
\[
        \sum_1^n u_i=0
\]
where $n$ is the number of nodes, should be build nodal point wise.

\subsection*{Distributed constraints}

As an example we take the periodic boundary condition Eq.~\eqref{eq:periodic}.
The condition Eq.~\eqref{eq:periodic} can be assumed to be valid exactly in the
nodal points of the curve (collocation):
\begin{equation}
      \vec\lambda_k\cdot\bigl( \vec u^k_1(y) - \vec u^k_2(y) \bigr)=0,
   \qquad\text{for all $\vec\lambda_k$, $k=1,\ldots,n$}
    \label{eq:collocationconnection}
\end{equation}
 or in an integral way (weak):
\begin{equation}
    \int_{C1} \vec\lambda\cdot\bigl( \vec u_1(y) - \vec u_2(y) \bigr)\D y=0,
   \qquad\text{for all $\vec\lambda$}
    \label{eq:weakconnection}
\end{equation}
The constraint are build nodal-point-wise and element-wise,
respectively\footnote{For a constraint in a point or on a nodeset it is only
possible to build it nodal-point wise, since elements are not defined on these
geometric entities.}.
In case of weak constraints on objects Eq.~\eqref{eq:weakconnection} is build
integration point wise, i.e.\ on element level the contribution to the integral
of a single integration point must be supplied.

It is possible to connect two parts of the domain
having elements that are not compatible
(mortars). 
One way to do this is using two curves (or two surfaces).
In that case, it is assumed that the Lagrange multipliers are
defined in the nodes of C1 and that all unknowns on C2 are involved in
all constraints (row dimension of $\mat A_{2e}$ is $\dim{\col u_2}$).
Another, more flexible, way to achieve this is to use connecting objects.

The heading of the element subroutine defining a constraint is given by
\begin{listing}{1}
   
  subroutine elemsub ( mesh, problem, constr, elem, node, matrix, vector, &
    first, last, coefficients, oldvectors, elemmat, elemmat2, elemmatadd, &
    elemvec, elemvecadd )
    type(mesh_t), intent(in) :: mesh
    type(problem_t), intent(in) :: problem
    integer, intent(in) :: constr, elem, node
    logical, intent(in) :: matrix, vector, first, last
    real(dp), intent(out), dimension(:,:) :: elemmat, elemmat2, elemmatadd
    real(dp), intent(out), dimension(:) :: elemvec, elemvecadd

     ....

  end subroutine elemsub

\end{listing}

with\\[2ex]
\begin{tabular}{ll}
\texttt{mesh, problem} & mesh and problem structure as in the heading of
                           \texttt{build\_system\_constraint}\\
\texttt{constr} & constraint number \\
\texttt{elem} & element number within this constraint \\
\texttt{node} & node or integration point number within this constraint\\
\texttt{matrix, vector} & logicals defined by \texttt{buildmatrix} and
                         \texttt{buildvector} in the heading \\
                        & of \texttt{build\_system\_constraint}.\\
\texttt{first, last} & logicals indicating whether the current element or node
    is the first/last \\ & in this constraint. \\
\texttt{coefficients, oldvectors} & coefficients and oldvectors structure as in the heading of
                           \texttt{build\_system\_constraint}\\
\texttt{elemmat} & matrix $\mat A$ or $\mat A_1$ for this element/node. \\
\texttt{elemmat2} & matrix $\mat A_2$ for this element/node. \\
\texttt{elemmatadd} & matrix $\mat B$ for this element/node. \\
\texttt{elemvec} & vector $\col g$ for this element/node. \\
\texttt{elemvecadd} & vector $\col h$ for this element/node. \\
\end{tabular}

\section{Non-linear constraints}
\label{sec:nonlinearconstraint}
Assume that we have a constraint minimization problem that is non-linear, i.e.\
the functional $I(\col u)$ is non-linear (not just quadratic) and the
constraint is non-linear as well:
\begin{equation}
   \col F(\col u) = \col 0
\end{equation}
For simplicity we have not included the additional unknowns.
The problem is equivalent to finding the stationary value of
\[
   J(\col u,\col\lambda)= I(\col u) +
    \col\lambda^T \col F(\col u )
\]
The stationary value is given by
\begin{equation}
   \delta {J} = \delta \col u^T(\pderiv{I(\col u)}{\col u}+
                                \pderiv{F(\col u)^T}{\col u}\col\lambda)+
                              \delta \col\lambda^T \col F(\col u ) = 0
\end{equation}
for all $\delta \col u$ and $\delta\col\lambda$, which leads to the system
\begin{align}
\pderiv{I(\col u)}{\col u}+
                 \pderiv{\col F(\col u)^T}{\col u}\col\lambda&=\col 0\\
\col F(\col u )\hspace*{2cm} & = \col 0
\end{align}
or written in index notation
\begin{align}
\pderiv{I}{u_i}+
                 \pderiv{\col F_l}{u_i}\lambda_l&= 0\\
\col F_k\hspace*{1.5cm} & = 0
\end{align}
Iteration is required for solving this nonlinear system of equations. For
example, a Picard type iteration scheme might lead to a nonsymmetric system
matrix. A full Newton-Raphson iteration scheme will lead to a symmetric Jacobian
$J$:
\begin{equation}
   J=\begin{pmatrix}
      \pderivxy{^2I}{u_i}{u_j}+\pderivxy{^2F_l}{u_i}{u_j}\lambda_l &
   \pderiv{F_m}{u_i} \\[3ex]
   \pderiv{F_k}{u_j} & 0
   \end{pmatrix}
\end{equation}

For the implementation of a non-symmetric Picard type iteration scheme, the 
element subroutine \texttt{elemsub1} in the call of
\texttt{build\_system\_constraint} can be used instead of the usual 
\texttt{elemsub}. With \texttt{elemsub1} the transposed matrices 
$\mat A^T$ or $\mat A_1^T$ (\texttt{elemmatt}), 
$\mat A_2^T$ (\texttt{elemmat2t}) and
$\mat B^T$ (\texttt{elemmataddt}) need to be filled by the user, giving the
possibility to build a non-symmetric matrix. In addition, element vectors added
to the right-hand side $\col f$ need to be filled (\texttt{elemvecf} and
\texttt{elemvecf2}).

For the implementation of the symmetric Newton-Raphson system the usual
\texttt{elemsub} element routine can be used to build the Lagrange multiplier
part. However the ``direct stiffness'' term 
$\plderivxy{^2F_l}{u_i}{u_j}\lambda_l$ and the corresponding right-hand side
need to be build using 
standard tools, such as build on an object using (\texttt{build\_system})
or \texttt{build\_system\_connection} if the constraint is two-sided, like in
fluid-solid interactions.

\section{Examples of problems with constraints}

\subsection{A Stokes problem on a unit square with periodic boundary
         conditions: channel problem with specified flowrate; collocation}
\label{sec:stokes2}
                                                                                
We consider the Stokes equation for a channel problem on a unit square
(see Fig.~\ref{fig:channel2}). On the upper and lower boundary the velocity is
set to zero. The velocities on curve C2 and C4 are assumed to be the same
(periodic).
\begin{figure}[ht!]
   \begin{center}
   \includegraphics{channel2}
   \end{center}
   \caption{Channel problem}
    \label{fig:channel2}
\end{figure}
                                                                                
The example file is \texttt{stokes2.f90} and can be obtained by typing at the
command line:
\begin{quote}
   \texttt{getexample stokes}
\end{quote}
                                                                                
The Stokes element routines
in the add-on `stokes' (see Sec.~\ref{sec:stokesaddon}) are used for
the building of the system. 

\begin{listing}{61}
! create mesh
                                                                                
  meshgen_options%elshape = 6
  meshgen_options%nx = nx
  meshgen_options%ny = ny
                                                                                
  call quadrilateral2d ( mesh, meshgen_options )
                                                                                
  call add_to_mesh ( mesh, curve=[-4] )     ! curve 5
                                                                                
\end{listing}
\begin{description}
   \item[line 69] here an extra curve (it becomes C5) is created that is
defined in the opposite direction as curve C4. The reason is that curves that
become connected using constraints must be defined in the same direction.
\end{description}
\begin{listing}{93}
  call define_constraint ( mesh, input_probdef, &
    physq=1, curve1=2, nglobalc=1 )
                                                                                
  call define_constraint ( mesh, input_probdef, &
    physq=1, curve1=2, curve2=5, discretization='collocation', exclude=3 )
                                                                                
  call problem_definition ( input_probdef, mesh, problem )
                                                                                
\end{listing}
\begin{description}
   \item[line 93--94] here the global constraint for the flow rate along curve
     C2 is defined. Note that only physical quantity 1 (velocity) is involved
     in the constraint.
   \item[line 96--97] here the constraint connecting curve C2 and C5 
     is defined. Note that only physical quantity 1 (velocity) is involved
     in the constraint. Also collocation is used to connect the curves. Note,
     that the first and the last point on the curves are excluded from the
     constraint. The reason is that in these points the velocities are imposed
     using essential conditions and thus completely removed from the equations.
     Imposing a constraint in these points leads to a singular matrix (row and
     column with only zeros).
\end{description}
\begin{listing}{115}
! create system matrix
                                                                                
  call create_sysmatrix_structure_base ( sysmatrix, mesh, problem )
  call create_sysmatrix_structure_constraint ( sysmatrix, mesh, problem )
  call finalize_sysmatrix_structure ( sysmatrix )
                                                                                
  call create_sysmatrix_data ( sysmatrix )
                                                                                
! build (assemble) matrix and vector from elements

  call build_system ( mesh, problem, sysmatrix, rhsd, elemsub=stokes_elem, &
    coefficients=coefficients )

  call build_system_constraint ( mesh, problem, sysmatrix, rhsd, &
    constraint1=1, elemsub=stokes_constr_flowr, addmatvec=.true., &
    coefficients=coefficients )

  call build_system_constraint ( mesh, problem, sysmatrix, rhsd, &
    constraint1=2, elemsub=stokes_constr_node_conn, addmatvec=.true., &
    coefficients=coefficients  )

\end{listing}
\begin{description}
   \item[line 117--121] here the system matrix is created. First the structure:
     base \fem\ structure (line 117) + constraints structure (line 118) and the
     finalizing of the structure (line 119).
     The data itself is allocated on line 121.
   \item[line 128--130] here the system matrix concerning the flowrate
     constraint is built using the element routine
     \texttt{stokes\_contr\_flowr}. Note that \texttt{addmatvec=.true.} to
     avoid the initializing of the matrix to zero. In the element routine
     \texttt{stokes\_contr\_flowr} the element
      matrix $\mat A_e$ is filled with
     \[
          \int_{\Gamma_e}\phi_i\D y
     \]
     and the elementvector is filled with the flow rate $Q$ in the first element
     along the curve. Note that the actual value of the flowrate has been
     specified in the \texttt{coefficients}.
   \item[line 132--134] here the system matrix concerning the periodic
     connection constraint is built using the element routine
     \texttt{stokes\_constr\_node\_conn}. Note that \texttt{addmatvec=.true.} to
     avoid the initializing of the matrix to zero. In the element routine
     \texttt{stokes\_constr\_node\_conn} the element
      matrix $\mat A_{1e}$ is filled with
     \[
          \begin{pmatrix}
             1 & 0 \\
             0 & 1
          \end{pmatrix}
     \]
     and the matrix $\mat A_{2e}$ is filled with
     \[
          \begin{pmatrix}
             -1 & 0 \\
             0 & -1
          \end{pmatrix}=-\mat A_{1e}
     \]
     The elementvector $\col g_e$ is set to zero.
\end{description}
\begin{listing}{148}

! write profile to a file for plotting with gnuplot

  call printtofile ( mesh, problem, 'vprofile.out', curve=2, &
    vector=velocity )
   
\end{listing}
Here the velocity profile is printed to a file.

Run the problem by typing
\begin{quote}
   \texttt{stokes2}
\end{quote}
Typing the command
\begin{quote}
  \texttt{viewfig}
\end{quote}
will show something like
Figs.~\ref{fig:velocity2}, \ref{fig:vorticity2}, 
      and \ref{fig:pressure2}. The velocity profile is plotted in
Fig.~\ref{fig:vprofile}.
\begin{figure}[ht!]
   \begin{center}
   \includegraphics[scale=0.5]{velocity2}
   \end{center}
   \caption{Plot of the velocity vector for the periodic channel problem}
    \label{fig:velocity2}
\end{figure}
\begin{figure}[ht!]
   \begin{center}
   \includegraphics[scale=0.5]{vorticity2}
   \end{center}
   \caption{Plot of the vorticity for the periodic channel problem}
    \label{fig:vorticity2}
\end{figure}
\begin{figure}[ht!]
   \begin{center}
   \includegraphics[scale=0.5]{pressure2}
   \end{center}
   \caption{Plot of the pressure for the periodic channel problem}
    \label{fig:pressure2}
\end{figure}
\begin{figure}[ht!]
   \begin{center}
   \includegraphics{vprofile}
   \end{center}
   \caption{Plot of the velocity profile in the periodic channel problem}
    \label{fig:vprofile}
\end{figure}

A similar example is available as \texttt{stokes33.f90}, but it is build using
nodesets instead of curves.

\subsection{A Stokes problem on a unit square with periodic boundary
         conditions: channel problem with specified flowrate; weak constraint}
\label{sec:weakflowrate}
                                                                                
Again we consider the Stokes equation for a channel problem on a unit square
(see Fig.~\ref{fig:channel2}), as in the previous section. The difference is
the weak implementation of the connection
between 
curves C2 and C4 where the velocities are assumed to be the same
(periodic) using Eq.~\eqref{eq:weakconnection}.

The example file is \texttt{stokes3.f90} and can be obtained by typing at the
command line:
\begin{quote}
   \texttt{getexample stokes}
\end{quote}
                                                                                
The Stokes element routines
in the add-on `stokes' (see Sec.~\ref{sec:stokesaddon}) are used for
the building of the system. 

Some notes on the source:
\begin{listing}{92}

  call define_constraint ( mesh, input_probdef, &
    physq=1, curve1=2, nglobalc=1 )
                                                                                
  call define_constraint ( mesh, input_probdef, &
    physq=1, curve1=2, curve2=5, discretization='weak', elementdof=[2,0,2] )
                                                                                
  call problem_definition ( input_probdef, mesh, problem )
                                                                                
\end{listing}
\begin{description}
   \item[line 96--97] here the constraint connecting curve C2 and C5 
     is defined. Note that only physical quantity 1 (velocity) is involved
     in the constraint. The connection is assumed to be weak and the Lagrange
     multipliers for both the $u$ and $v$ are interpolated using linear shape
     functions. This is reflected in the parameter
     \texttt{elementdof=[2,0,2]}, which represents the degrees of freedom of
     the Lagrange multipliers in the nodes of an element along the curve.
     In the weak implementation it is not necessary to exclude any points
     from the constraint, in contrast to the collocation implementation.
\end{description}
\begin{listing}{123}

! build (assemble) matrix and vector from elements

  call build_system ( mesh, problem, sysmatrix, rhsd, elemsub=stokes_elem, &
    coefficients=coefficients )

  call build_system_constraint ( mesh, problem, sysmatrix, rhsd, &
    constraint1=1, elemsub=stokes_constr_flowr, addmatvec=.true., &
     coefficients=coefficients )

  call build_system_constraint ( mesh, problem, sysmatrix, rhsd, &
    constraint1=2, elemsub=stokes_constr_elem_conn, addmatvec=.true., &
     coefficients=coefficients )

\end{listing}
\begin{description}
   \item[line 133--135] here the system matrix concerning the periodic
     connection constraint is built using the element routine
     \texttt{stokes\_constr\_elem\_conn}. Note that \texttt{addmatvec=.true.} to
     avoid the initializing of the matrix to zero. In the element routine
     \texttt{stokes\_contr\_elem\_conn} the element
      matrix $\mat A_{1e}$ is filled as follows
\[
   \mat A_{1e}=
   \begin{pmatrix}
       \int_{e}\psi_i\phi_j\D y & 0 \\
       0 & \int_{e}\psi_i\phi_j\D y 
   \end{pmatrix}
\]
     where $\psi_i$ are the linear shape functions
     of $\lambda$ and $\psi_j$ the quadratic shape functions of the velocity.
The matrix contains four partitions.
The row partitioning represents the
Lagrange multiplier in horizontal ($x$) and vertical ($y$) direction,
respectively. The
column partitioning represents the contribution of the $u$ velocity component
and the $v$ velocity component to the constraints, respectively.
The contribution of the $v$ ($u$) component to the horizontal (vertical)
direction is zero \footnote{In \tfem\ it is assumed that \emph{all}
 components of 
a physical quantity contribute to \emph{all}
constraint directions. Therefore the off-diagonal
blocks have to be filled with zeros here. The only way to avoid this is to
define each velocity component as a separate physical quantity.}.
The matrix $\mat A_{2e}$ is equal to $-\mat A_{1e}$ and
     the elementvector $\col g_e$ is set to zero.
\end{description}
Run the problem by typing
\begin{quote}
   \texttt{stokes3}
\end{quote}
and view the plot files with
\begin{quote}
  \texttt{viewfig}
\end{quote}

\subsection{A Stokes problem on a unit square with periodic boundary
         conditions: channel problem with specified pressure difference using
         additional unknowns; weak constraint}
                                                                                
In the source code \texttt{stokes4.f90} the flow rate $Q$ is defined as an
additional unknown:
\[
     \int_{C1}u_x\D y - Q = 0
\]
and a matrix $\mat B$ is filled. The single component of the
right-hand side vector $\col h$
is now $\Delta p$, the pressure difference between the inflow
and outflow of the channel. This is a bit of an artificial problem, since the
pressure difference can easily be imposed using natural boundary elements on
either C2 or C4 without using any constraints.

\subsection{A Stokes problem on a unit square with periodic boundary
         conditions: channel problem with specified flowrate; collocation on
         two objects}
                                                                                
We consider again the Stokes equation for a channel problem on a unit square
(see Fig.~\ref{fig:channel2}). On the upper and lower boundary the velocity is
set to zero. The velocities on curve C2 and C4 are assumed to be the same
(periodic). In this problem we connect the flow exit and entry using two 
objects.
                                                                                
The example file is \texttt{stokes15.f90} and can be obtained by typing at the
command line:
\begin{quote}
   \texttt{getexample stokes}
\end{quote}
                                                                                
The Stokes element routines
in the add-on `stokes' (see Sec.~\ref{sec:stokesaddon}) are used for
the building of the system. 

\begin{listing}{132}
                                                                                
   call quadrilateral2d ( mesh, meshgen_options )

   call add_to_mesh ( mesh, object='curve', objectcurve=2, &
     excludepoints=[2,3] )
   call add_to_mesh ( mesh, object='curve', objectcurve=2, &
     excludepoints=[2,3] )

 ! shift object 2 to entry side
   mesh%objects(2)%coor(:,1) = mesh%objects(2)%coor(:,1) - 1._dp

   call fill_mesh_parts ( mesh )

\end{listing}
\begin{description}
   \item[line 135--136] here an object is created on curve C2, using the curve
C2. In fact, the object has the same coordinates as the curve except for the
first and last point of the curve (due to \texttt{excludepoints=[2,3]}).
The first and the last point of the curve has been excluded
     because in these points the velocities are imposed
     using essential conditions and thus completely removed from the equations.
     Imposing a constraint in these points leads to a singular matrix (row and
     column with only zeros).
   \item[line 137--138] here again an object is created on curve C2,
using the curve C2. In fact, the object is identical to the previous object.
   \item[line 141] here the coordinates of the second object are shifted
to the entry section.
\end{description}
\begin{listing}{167}

   call define_constraint ( mesh, input_probdef, &
     physq=1, object=1, object2=2, discretization='collocation', nodedof=2 )

\end{listing}
\begin{description}
   \item[line 168--169] here the constraint connecting the exit and entry
section is defined using the previously defined objects.
     Note that only physical quantity 1 (velocity) is involved
     in the constraint. Also collocation is used to connect the objects.
\end{description}
\begin{listing}{203}
                                                                                
   call build_system_constraint ( mesh, problem, sysmatrix, rhsd, &
     constraint1=2, elemsub=constraints_object_conn, addmatvec=.true. )

\end{listing}
\begin{description}
   \item[line 204--205] here the system matrix concerning the periodic
     connection constraint is built using the element routine
     \texttt{constraints\_object\_conn}. Note that \texttt{addmatvec=.true.} to
     avoid the initializing of the matrix to zero. In the element routine
     \texttt{constraints\_object\_conn} the element
      matrix $\mat A_{1e}$ and  $\mat A_{2e}$ are filled with
     \[
   \mat A_{e1}=
          \begin{pmatrix}
             \phi_{1i}(\vek x_{c1}) & 0 \\
             0 & \phi_{1i}(\vek x_{c1})
          \end{pmatrix}
     \]
and 
     \[
   \mat A_{e2}=
          \begin{pmatrix}
             -\phi_{2i}(\vek x_{c2}) & 0 \\
             0 & -\phi_{2i}(\vek x_{c2})
          \end{pmatrix}
     \]
where $\vek x_{c1}$, $\vek x_{c2}$ are the
position vectors of the current collocation point in both objects and 
 $\phi_{1i}$, $\phi_{2i}$ are the shape functions of the velocity in
the corresponding elements where the connection points are.
\end{description}

REMARK: the technique described in this section has essential the same result
as the technique in Sec.~\ref{sec:stokes2} using connecting curves. However,
the current technique of connecting objects is more general:
\begin{itemize}
   \item curves, surfaces need to have exactly the same topology to get
connected. Especially in 3D this is a big problem for complicated meshes.
Objects don't suffer from that problem, because they connect
`volume points', not boundary points.
   \item since objects connect `volume points', it is also possible to define
constraints using the gradients of the solution.
   \item points in the objects don't have to be in nodal points. This means
that curves/surfaces can be connected having different element distribution or
resolution. One should be careful here, though, since collocation might not
always be a mathematical sound way to connect boundaries having different
resolution. 
\end{itemize}

\subsection{A Stokes problem on a unit square with periodic boundary
         conditions: channel problem with specified flowrate; weak constraint;
         global constraint on an object}

Again we consider the Stokes equation for a channel problem on a unit square
(see Fig.~\ref{fig:channel2}), as in section~\ref{sec:weakflowrate}.
The difference is that the flow rate constraint is not on an end curve but on
a cross section that can be positioned anywhere along the channel.
The global constraint for the flow rate is on an object that is made from a
curve and shifted to any position. The pressure jump is now also across the
object. This example is somewhat artificial, but is shows the possibilities
of global constraints on objects.

The example file is \texttt{stokes18.f90} and can be obtained by typing at the
command line:
\begin{quote}
   \texttt{getexample stokes}
\end{quote}
The example has been made such that it is easy to switch between continuous
and discontinuous pressure interpolation.

Some notes:
\begin{itemize}
   \item For discontinuous pressure ($P_1$) elements the pressure jump
(=Lagrange multiplier) seems to be exact, even if the object is positioned
inside elements. The same is true for the velocity profile.
   \item For continuous pressure ($Q_1$) elements the pressure jump
(=Lagrange multiplier) is lower than the exact value. The lowest value
is obtained if the object is positioned at the interface between elements (in
the example the error is 4\%) and
the best value is if the object is half way an element (error 0.4\%).
Correspondingly the
velocity profile in the channel is proportionally lower than the exact one.
This means that in the elements at the pressure jump near the object
``some fluid volume is lost''. This is related to fact that
Taylor-Hood elements are
only globally volume conserving (``incompressible'')
and not locally element by element.
The error improves with mesh refinement, but only linearly.
\end{itemize}
The above notes suggest that discontinuous elements are better for pressure
jumps, like in fictitious domain methods. In practice, however,
the differences are likely to be small with sufficient refinement, but it is
wise to test both type of elements for your problem. Furthermore, discontinuous
pressure elements lead to instabilities in time-dependent viscoelastic flows
and are therefore not really an option in that case.
